\extrachap{Acronyms and Nomenclature}

\section*{Acronyms}
\begin{description}[\settowidth{\labelwidth}{\textbf{OpenRadioss}}]
\item[\textbf{1U}] One-Unit CubeSat standard form factor ($100 \times 100 \times 113.5\text{ mm}$)
\item[\textbf{AlSi10Mg}] Aerospace-grade aluminum-silicon-magnesium casting alloy for L-PBF
\item[\textbf{B-Rep}] Boundary Representation (solid CAD geometry)
\item[\textbf{CAD}] Computer-Aided Design
\item[\textbf{CDS}] CubeSat Design Specification (CalPoly standard)
\item[\textbf{DfAM}] Design for Additive Manufacturing
\item[\textbf{DoE}] Design of Experiments
\item[\textbf{FEA}] Finite Element Analysis
\item[\textbf{FEM}] Finite Element Method
\item[\textbf{FGL}] Functionally Graded Lattice
\item[\textbf{FS}] Factor of Safety
\item[\textbf{GEVS}] General Environmental Verification Standard (NASA GSFC-STD-7000A)
\item[\textbf{LEFM}] Linear Elastic Fracture Mechanics
\item[\textbf{LEO}] Low Earth Orbit
\item[\textbf{L-PBF}] Laser Powder Bed Fusion (metallic 3D printing)
\item[\textbf{MS}] Margin of Safety
\item[\textbf{NASA}] National Aeronautics and Space Administration
\item[\textbf{NSGA-II}] Non-dominated Sorting Genetic Algorithm II
\item[\textbf{NSIF}] Notch Stress Intensity Factor
\item[\textbf{P-POD}] Poly-Picosatellite Orbital Deployer
\item[\textbf{PGML}] Physics-Guided Machine Learning
\item[\textbf{PSD}] Power Spectral Density ($\text{g}^2/\text{Hz}$)
\item[\textbf{Q}] Quality Factor ($Q = 1 / (2\zeta)$)
\item[\textbf{ResNet}] Residual Neural Network
\item[\textbf{RMS}] Root Mean Square
\item[\textbf{SRS}] Shock Response Spectrum
\item[\textbf{Tet10}] 10-node quadratic tetrahedral finite element
\item[\textbf{TOPSIS}] Technique for Order of Preference by Similarity to Ideal Solution
\item[\textbf{TRL}] Technology Readiness Level
\end{description}

\section*{Key Mathematical Symbols}
\begin{description}[\settowidth{\labelwidth}{$a_c, a_{\text{init}}$}]
\item[$a_c, a_{\text{init}}$] Critical and initial crack sizes in linear elastic fracture mechanics
\item[$D$] Palmgren-Miner cumulative fatigue damage index ($D \le 0.25$)
\item[$E, E^*$] Solid and effective cellular Young's modulus of elasticity
\item[$f_1$] First fundamental structural eigenfrequency ($\text{Hz}$)
\item[$FS_{\text{yield}}, FS_{\text{ult}}$] NASA qualification safety factors for yield ($1.25$) and ultimate ($1.40$)
\item[$K_t, K_I$] Elastic stress concentration factor and Mode-I stress intensity factor
\item[$\Delta K_{\text{th}}$] Fatigue crack propagation threshold ($\text{MPa}\sqrt{\text{m}}$)
\item[$m$] Basquin fatigue S-N curve slope exponent
\item[$MS$] Structural Margin of Safety ($MS = \sigma_{\text{adm}}/\sigma_{3\sigma} - 1.0$)
\item[$\zeta$] Critical viscous damping ratio ($\zeta = 0.02$, $2\%$)
\item[$\theta$] Re-entrant auxetic internal cell angle
\item[$\nu, \nu_{\text{eff}}$] Solid Poisson's ratio and effective lattice Poisson's ratio
\item[$\rho_s, \rho^*$] Base material density and effective cellular density
\item[$\sigma_{3\sigma}$] Peak stochastic stress at three-sigma confidence level ($99.73\%$)
\item[$\sigma_{\text{yield}}, \sigma_{\text{ult}}$] Material yield strength and ultimate tensile strength
\item[$\sigma_{\text{adm}}$] Admissible design stress ($\sigma_{\text{yield}} / FS_{\text{yield}}$)
\end{description}
